% =============================================================================
% SUPPLEMENT S5 -- Robustness: conformal secondary scores and the
% pre-registered sensitivity slices.
%
% Carries Appendix A of the full study (paper/main.tex) into the supplement.
% Numbering (S-sections, S-tables, S-figures) is set by supplement.tex; this
% file deliberately introduces no table or figure float, because the
% manuscript cites the infeasible-cell ledger of Section S6 as Tables S11 and
% S12 and that numbering depends on every float printed before it.
% =============================================================================

\section{Robustness: conformal secondary scores and the pre-registered
  sensitivity slices}
\label{supp:robustness}

This section carries the two robustness items the manuscript defers to the
supplement. Section~\ref{supp:robustness-conformal} states the scoring
convention that makes CRPS, the log score and PIT \emph{placeholders} for
every conformal arm: it is the convention behind the manuscript's rule in
Section~4 that those values are flagged \texttt{secondary}, read only against
other conformal cells and never ranked against distributional mechanisms;
behind the sentence in Section~5 that conformal PIT values are placeholders;
and behind the limitation in Section~8 that comparing a conformal cell's CRPS
with a native cell's compares a convention with a forecast.
Section~\ref{supp:robustness-slices} documents the pre-registered secondary
runs -- the sensitivity slices announced in Sections~3 and~8 of the manuscript
and registered in Addendum~1 -- says exactly what the released file carries for
each of them, and Section~\ref{supp:robustness-limits} states plainly the two
things those slices cannot do. Neither item changes a headline number: the
conformal convention governs which scores may be compared with which, and the
slices re-score the same fits on pre-specified subsets of rows.

\subsection{Proper scores for conformal cells}
\label{supp:robustness-conformal}

Conformal wrappers emit intervals, not distributions. The harness preserves a
single sampling interface across all seven uncertainty mechanisms by drawing
uniformly inside each interval, so CRPS, the log score and the PIT computed
from those draws describe a placeholder density rather than the wrapper's own
predictive law. They are carried per cell in the released results files under a
\texttt{secondary} flag, are read only against other conformal cells, and are
never ranked against distributional mechanisms (manuscript Section~4). They are
deliberately not printed beside the distributional scores of the manuscript's
Section~6 or of Section~S3 here: a reader who compares a conformal cell's CRPS
with a native cell's CRPS is comparing a convention with a forecast, and a
shared table invites exactly that. The same convention governs the derived
quantities, whose $e_0$, $e_{65}$ and annuity quantiles for a conformal cell
come from the filler draws and are therefore reported as not applicable
throughout (Addendum~2, \S3); those are the \texttt{n/a} rows of the full
actuarial table in Section~S3.

What the convention does \emph{not} touch is any conformal number that enters a
reported contrast. Coverage, interval width, the interval score and joint path
coverage for a conformal cell are read from the wrapper's own bounds at its own
construction level of 95\%, with no Poisson composition -- the radius already
contains observation noise -- and no quantile-of-samples step (Addendum~2~\S3,
Addendum~3~\S6). Coverage and interval scores at the 50\% and 80\% levels are
not applicable to conformal cells and are reported as such, and any contrast
involving a conformal arm is decided on the 95\% interval score. The placeholder
scores exist only so that one code path can score all fifty cells; they are
released for completeness and for anyone wishing to re-derive them, not to be
ranked.

\subsection{The pre-registered sensitivity slices}
\label{supp:robustness-slices}

The pre-registered secondary runs re-score the headline coverage contrasts on
pre-specified subsets of the same fits rather than refitting anything. No model
is estimated a second time, no seed changes and no split moves; a slice selects
rows of the released per-cell results and recomputes the same means over them.
Each slice is computed by \texttt{scripts/sensitivities.py} into
\texttt{results/sensitivities.json}, released with the code.

\paragraph{Shift-regime slices.} The baseline \texttt{full\_panel} retains all
twenty shift populations (40 population--sex units, 2{,}000 cell rows, 210 of
them error rows). Against that baseline:
\begin{itemize}
\item \texttt{drop\_2024} -- the provisional-revision sensitivity. The panel is
  unchanged; the test window is restricted to 2020--2023 and coverage is the
  mean of the retained per-horizon columns, $h=1,\dots,4$.
\item \texttt{drop\_USA\_CHL} -- the two populations whose exposures HMD
  documents as census-based and subject to later revision are removed, leaving
  eighteen populations (36 units, 1{,}800 rows, 180 error rows;
  Addendum~1~\S B.1).
\item \texttt{register\_based} -- CHE, DNK, FIN, ISL, NOR and SWE: six
  populations (12 units, 600 rows, 8 error rows).
\item \texttt{census\_based} -- the remaining fourteen populations (28 units,
  1{,}400 rows, 202 error rows).
\end{itemize}
The register-versus-census split as computed follows
\texttt{docs/DATA-PREREQS.md}~\S B.2 item~9, which places EST, LVA and LTU in
the census-based set because their 2021--2024 exposures rest on post-censal
inputs. Addendum~1~\S B's prose lists those three as register-based in passing,
and Section~3 of the manuscript repeats that description. The released file uses
the \texttt{DATA-PREREQS} definition and records the discrepancy in its own
\texttt{note} field rather than resolving it silently; it is restated here for
the same reason. A reader who prefers the other definition can recompute both
slices from the released per-cell rows.

\paragraph{Placebo-regime slices.} The baseline \texttt{full\_panel} retains all
eleven placebo populations (22 units, 1{,}100 rows, 136 error rows). Against
that baseline:
\begin{itemize}
\item \texttt{drop\_GBR\_SCO} -- GBR\_SCO, whose 1912--1938 population is
  interpolated by a spline HMD itself calls unreliable, is removed, leaving ten
  populations (20 units, 1{,}000 rows, 128 error rows).
\item \texttt{neutral\_only} -- the neutral stratum CHE, DNK, FIN, ISL, NLD,
  NOR and SWE, the cleanest pandemic-only comparison: seven populations
  (14 units, 700 rows, 90 error rows).
\item \texttt{dnk\_full\_window} and \texttt{dnk\_drop\_1921\_1922} -- Denmark
  alone (2 units, 100 rows, 8 error rows), scored first over the full
  1914--1922 window and then over 1914--1920 only, Denmark having gained South
  Jutland in 1921, inside the test window.
\end{itemize}
The three pre-specified placebo strata of Addendum~1~\S A are neutral (CHE,
DNK, FIN, ISL, NLD, NOR, SWE), belligerent-total (FRATNP, GBRTENW, ITA) and
civilian-only (GBR\_SCO). Only the neutral stratum is a slice of its own here;
the stratum-level coverage that the manuscript's Section~6 reads is in the
twin-crises table of Section~S3, and the pooled placebo result is always
reported alongside the strata.

\paragraph{What the file carries.} For every family--mechanism arm in every
slice -- all fifty primary cells, whether or not the arm survives the slice --
the file records mean 95\% coverage over the rows the slice retains, joint path
coverage wherever the slice keeps the whole path, the number of scored cells,
the number of error rows behind that mean, the number of ages scored, and flags
for whether the arm is conformal and whether it is a secondary grid cell. Each
slice also carries a metadata block giving the retained populations, the row and
cell counts quoted above, and the error rows decomposed into the quality gate's
classes (machine, design-floor, method and residual), so that no mean in the
file has a denominator a reader cannot reconstruct. Means are unweighted over
valid rows, one row per (population, sex, origin) cell and arm, and arms are
never pooled. The same file carries two further blocks the manuscript relies on:
the per-origin effective cluster counts behind the clustered inference of
Section~5 -- at the 2019 shift origin, for instance, twenty populations
contribute at least one valid row but only seven contribute a valid row to every
one of the fifty primary arms -- and the common-cell loss panels on which the
Diebold--Mariano and model-confidence-set contrasts of Section~6 are computed.

\subsection{Two limits, stated plainly}
\label{supp:robustness-limits}

\paragraph{A slice that drops a horizon cannot report joint path coverage.} The
runner emits one indicator for the whole forecast path rather than the
per-horizon indicators a shortened path would need, so joint path coverage is
not recoverable for the two horizon-restricted slices,
\texttt{shift/drop\_2024} and \texttt{placebo/dnk\_drop\_1921\_1922}. Their
\texttt{joint\_path\_coverage\_95} entries are null rather than imputed. What
the file does record for them is
\texttt{joint\_path\_coverage\_95\_lower\_bound}, the full-path value: a path
that covers at every horizon $1,\dots,H$ also covers at every horizon of a
subset of them, so the full-path rate is a valid lower bound for the shortened
path and is labelled as such rather than reported as the quantity itself. No
joint-coverage claim anywhere in the manuscript rests on a horizon-restricted
slice.

\paragraph{The registered age-cap sensitivity is not among the slices.} The
pre-registration names a run with ages capped at 90 rather than 99, to test
whether the old-age reading of \Hfour\ is an artefact of noisy old-age cells in
small populations. That run is not among the slices in
\texttt{results/sensitivities.json}: recomputing it needs per-age scores, which
the released per-cell rows do not carry, so it would require a re-score rather
than a re-selection of rows. What this paper reports about the old-age stratum
instead is the age-band decomposition of the manuscript's Section~6, which
scores $0$--$24$, $25$--$64$ and $65$--$99$ separately -- the same partition the
conformal wrappers are Mondrian-stratified on -- with the full per-cell profile
in Table~\ref{tab:h4-age} of Section~S3 and the single-age coverage curves in
the coverage-by-age figures of Section~S4. The age-cap run itself remains
outstanding, and is recorded here as outstanding rather than treated as
discharged by the age-band decomposition.
